\section*{Введение}
\addcontentsline{toc}{section}{Введение}

\textbf{Актуальность темы.} Библиотеки хранят и обрабатывают значительный объем
информации о книжном фонде, читателях, выдаче и возврате книг. При ручном ведении
учета или использовании разрозненных таблиц возникают проблемы с поиском данных,
контролем доступности книг, учетом задолженностей и формированием статистики.
Разработка веб-приложения для автоматизации учета книг и читателей позволяет
упростить работу библиотекаря и сделать информацию о фонде более доступной для
читателей.

\textbf{Объект исследования} --- процессы учета книг, читателей и выдачи книг в
библиотеке.

\textbf{Предмет исследования} --- веб-приложение для автоматизации учета книг и
читателей библиотеки.

\textbf{Цель работы} --- разработать веб-приложение для автоматизации учета книг
и читателей библиотеки.

Для достижения поставленной цели необходимо решить следующие задачи:

\begin{enumerate}
    \item Выполнить анализ предметной области и существующих программных решений.
    \item Выбрать программные инструменты разработки веб-приложения.
    \item Спроектировать структуру приложения и модель данных.
    \item Реализовать аутентификацию, авторизацию и разграничение прав доступа.
    \item Реализовать CRUD-интерфейс для книг, читателей и журнала выдач.
    \item Реализовать поиск, фильтрацию, статистику, загрузку и выгрузку файлов.
    \item Подготовить отчет о проделанной работе.
\end{enumerate}

\textbf{Практическая значимость работы} заключается в создании учебного
веб-приложения, демонстрирующего основные принципы разработки информационной
системы: работу с пользователями, базой данных, интерфейсом, файлами и правами
доступа.